\section{The Foundation}
\label{sec:foundation}

The model has four named quantities. Everything in this section follows from the axiom and the two constraints of Section~\ref{sec:constraints}.

\begin{definition}[Vibe]
A \emph{vibe} is an indivisible unit of experience. It is the only kind of entity. We write $V$ for the population of vibes.
\end{definition}

\begin{definition}[Tone]
Each vibe $v$ carries an intrinsic experiential quality, its \emph{tone} $t_v$, drawn from a finite alphabet $T$. The working choice is ternary,
\[
t_v \in T = \{-1, 0, +1\},
\]
read as pain, peace, and pleasure. The value $0$, peace, is the absence of a committed pole, a vibe whose influence is too weak to settle either way. Binary $\{-1,+1\}$ is the strict minimum, recovered by forbidding $0$. Richer finite alphabets are reserved for later (Section~\ref{sec:matter}). Continuous tone and complex phase are emergent only.
\end{definition}

\begin{definition}[Note]
The \emph{note} is the only primitive relation. To say $v$ notes $w$ is to say that $v$ experiences, or is present to, $w$. The set of vibes a vibe notes is its \emph{neighborhood} $N(v)$.
\end{definition}

Rich experience is meant to come from patterns of simple tones across many vibes, not from a complicated tone inside one vibe. A single ternary vibe knows only three states. Joy, sorrow, love, and fear are configurations of vast meshes, the way binary digits make unlimited information by combination.

\subsection{What a link is}

The principle of monism forbids a separate link structure. So a ``link'' between two vibes is never a third thing between them. It is one of two cases.

\begin{definition}[Bare note and relational vibe]
A link between $a$ and $b$ is either a \emph{bare note}, the simple fact that $a$ notes $b$ and $b$ notes $a$, carrying no stored data, or a \emph{relational vibe}, a vibe whose neighbors are the two it relates and whose tone is the quality of the relation. More elaborate links are small structures of such vibes.
\end{definition}

No infinite regress follows. The bare note is unmediated. A vibe and a relational vibe note each other directly, with nothing between. Relational vibes appear only where a relation needs its own state. This is forced by the axiom. A link that existed as a thing would have to be experiential, hence a vibe, so anything real that a link carries is already a vibe.

\subsection{Will and fill, the two directions of a note}

The axiom's two clauses, read along a note, are the model's two directions of influence.

\begin{definition}[Will]
The \emph{will} of a vibe is its tone turned outward, $\mathrm{will}_v = \omega(t_v)$, the part of it that another vibe reads when noting it. Will is the outside of a tone, as the felt quality is its inside. They are one thing seen from two sides. Nothing is stored on a link.
\end{definition}

\begin{definition}[Fill]
The \emph{fill} of a vibe toward a neighbor is how strongly it takes that neighbor in. It is carried by the note itself, the strength of attention being the tone of the relational vibe between the two.
\end{definition}

\begin{proposition}[Push and Pull]
Every note carries influence in both directions at once, active projection (will) and receptive integration (fill). Each vibe simultaneously influences and is influenced by the vibes it notes.
\end{proposition}

Base will carries no purpose. It is the unavoidable outward presence of a vibe that, by the axiom, can be experienced. Volition, the directed wanting of an agent, is not primitive. It emerges only for coherent meshes (Section~\ref{sec:mind}). The word ``will'' names the push, not an agent.

\subsection{The vibe mesh}

\begin{definition}[Vibe Mesh]
The \emph{vibe mesh} is the whole relational system, the population of vibes, their tones, and the note relation among them. The same word names the global totality or any coherent sub-structure within it. A mind is a vibe mesh, as is a cell, an atom, or any self.
\end{definition}

There is no external space or time holding the mesh. All structure is internal to it, and space and time are recovered from the note-structure later (Sections~\ref{sec:time} and~\ref{sec:geometry}).

\subsection{A concrete minimal case}

The smallest non-trivial mesh is two vibes $a, b$ that note each other, like two coupled Ising spins or two adjacent cells in a cellular automaton. Each carries a tone in $\{-1,0,+1\}$. The note $a \to b$ is the channel by which $a$ reads $b$'s tone. The will of $a$ is the tone it broadcasts, $\omega(t_a)$. Its fill is how strongly it weighs $b$, carried by the strength of the note between them. If that note acquires its own state, a coupling that strengthens or weakens, it is a third vibe, a relational vibe. Nothing else is present. Two tones, the note, and at most the note's own tone.
